\section{从一个纽结多项式说起}

\subsection{Jones 多项式的定义}

设 $L\subset S^3$ 是一个定向链环，画成平面投影图。在每个交叉处，我们可以做三种局部操作：
把交叉换成两条平行线（0-光滑化）、换成另一种平行线（1-光滑化），或者保留交叉。
Jones 多项式由拆接关系（skein relation）刻画：
\begin{equation}
t^{-1}J(L_+)-t\,J(L_-)=(t^{1/2}-t^{-1/2})\,J(L_0), \tag{1.1}
\end{equation}
其中 $L_+,L_-,L_0$ 是三个投影图，除了一个交叉外完全相同，该交叉分别是正交叉、负交叉与光滑化。
再要求平凡纽结的值为 $1$，就唯一确定了链环的多项式 $J(L)\in\ZZ[t^{\pm1/2}]$。
例如三叶结的 Jones 多项式是
\[
J(\text{三叶结})=-t^{-4}+t^{-3}+t^{-1}.
\]

这个定义有几个漂亮的性质：它给出一个容易计算的算法；它把纽结的三种基本变换
（Reidemeister 变换）都吸收进了多项式；它还自然地带上一个变量 $t$，可以取任意复数。
但它也有一个根本的缺点：\textbf{定义依赖投影图}。为了证明 $J$ 是链环的不变量，
必须验证它在 Reidemeister 变换下不变；而在计算过程中，中间多项式来自完全不同的投影，
彼此之间没有任何内蕴的联系。

这正是 1980 年代纽结理论的核心困惑。人们知道 Jones 多项式存在、知道它很好用，
但不知道它“是什么”：它是不是某个更基本的、定义在三维空间里的量的影子？

\subsection{一个类比：1930 年代的量子电动力学}

Witten 在论文开头用了一个类比\Wp{352}。1930 年代，物理学家已经知道量子电动力学是正确的：
它给出与实验吻合的能级与散射截面。但当时没有人知道如何把 Maxwell 方程与量子力学
写成一种“明显相对论性”的形式；直到量子场论（特别是路径积分与协变微扰论）建立之后，
这个表述才出现。

Jones 多项式的处境类似：它的组合定义是“二维的”，而它最终应该来自某个“三维的”结构。
用 Witten 的话说，问题是：二维共形场论里究竟是哪个东西是三维的？

\subsection{二维的来源太多}

1980 年代中后期，纽结多项式从至少三条互不相干的路线冒出来\Wp{352}：
\begin{enumerate}[leftmargin=*]
  \item 可解格点模型与统计力学。某些二维格点模型的配分函数在临界点附近给出纽结不变量；
        相关的组合对象是 Temperley--Lieb 代数与 Yang--Baxter 方程。
  \item 二维共形场论的单值性。在有理共形场论里，共形块的单值性给出辫群表示，
        而辫群的迹与纽结多项式有关。
  \item 量子群与 Hopf 代数。Drinfeld 与 Jimbo 的量子群给出 Yang--Baxter 方程的解，
        进而给出纽结不变量。
\end{enumerate}
这三条路线各自都能算出 $J$，但它们彼此之间看不到明确的联系。一个自然的愿望是：
\textbf{为这种多样性找到一个统一的主题}。

\subsection{不变量家族与量子群}

Jones 多项式不是孤立的现象。1980 年代中后期，人们发现了一整族“量子不变量”：
HOMFLY 多项式（两个变量，来自 $SU(N)$ 的 $N$ 趋于一般值）、Kauffman 多项式（来自 $SO(N)$）、
以及更一般的 Reshetikhin--Turaev 不变量（来自量子群与模张量范畴）。
这些不变量有共同的结构：一个 $R$ 矩阵（Yang--Baxter 方程的解），配上辫群的表示，
再取迹。

这个观察很重要，因为它提示：纽结不变量不是某个二维组合技巧的产物，
而是某种“表示论 + 迹”的普遍结构。Witten 的工作给这个结构一个三维解释：
$R$ 矩阵来自 Wilson 线的融合，迹来自三维理论在曲面上的态空间。于是“量子群”
不再是代数上的巧合，而是三维拓扑场论的二维投影。

理解这一点，就能理解为什么后面要花大篇幅讲仿射李代数、共形块与融合规则：
这些看起来是二维的工具，实际上是在描述三维理论的态空间。

\subsection{一个不依赖度规的量子场论}

如果要用三维量子场论来定义纽结不变量，这个理论必须有一个非常特殊的性质：
它的配分函数不能依赖 $S^3$ 的度规，否则把纽结放进去时，结果会随度规改变。
换句话说，我们需要的是一种“拓扑场论”：定义在任意定向三维流形上、只依赖拓扑的量子场论。

这一点在广义相对论里有先例：引力是“对度规积分”的理论，它的定义不预设背景度规。
但广义相对论不是唯一的可能性。事实上，只要拉格朗日量本身不含度规，
量子化之后的配分函数就很可能不依赖度规。Witten 的观察是：三维的 Chern--Simons 理论
正是这样的理论，而它的 Wilson 线恰好给出纽结不变量\Wp{351}。

于是故事的目标变得清楚了：
\begin{enumerate}[leftmargin=*]
  \item 写出 Chern--Simons 作用量（第二讲）；
  \item 证明它不依赖度规、且耦合常数必须取整数（第二讲）；
  \item 把它量子化，得到曲面上的 Hilbert 空间（第三、四讲）；
  \item 把纽结写成 Wilson 线，得到一个三维不变量（第五讲）；
  \item 证明这个不变量就是 Jones 多项式，并算出来（第六、七讲）。
\end{enumerate}

这五步就是本讲义的主线。每一步都会遇到一个具体的技术问题：作用量如何变分、相空间是什么、
量子化用什么极化、穿刺点上的约束如何解、手术如何算。我们依次处理。

\subsection*{练习}
\begin{enumerate}[leftmargin=*]
  \item 用拆接关系 (1.1) 计算三叶结的 Jones 多项式，并与正文结果核对。
  \item 证明：若 $L$ 是平凡链环的 $n$ 个分量，则 $J(L)=(-t^{1/2}-t^{-1/2})^{n-1}$。
  \item 说明为什么 Jones 多项式不是链环的完全分类（给出两个不同纽结但同多项式的例子，或说明为什么难找）。
\end{enumerate}

\begin{quote}
\small
\textbf{本讲要点。} Jones 多项式有组合定义，但缺少内蕴的三维解释。
三维拓扑场论的配分函数（不依赖度规）与 Wilson 线的期望值，是可能的三维来源。
\end{quote}
